% ============================================================
% Round 06: GRPO Improvements
% ============================================================

\subsection{Round 06：GRPO 的改进方向}

\subsubsection{本轮问题}

\begin{conceptbox}
\textbf{学习目标：} 在掌握 GRPO 基础之后，理解后续改进主要在修什么问题：长度偏置、无效 group、clip 过强、token-level ratio 高方差、局部归一化偏置、过度思考。\\
\textbf{主线位置：}
\[
\text{GRPO}
\to
\text{DAPO / Dr.GRPO / GSPO / REINFORCE++ / DCPO / TreeRPO}
\]
\end{conceptbox}

\subsubsection{为什么还要改进 GRPO}

GRPO 的核心优点是：

\[
\text{去掉 critic，用 group relative reward 估计 advantage}
\]

但标准 GRPO 在大规模 reasoning RL 里会遇到几个实际问题：

\begin{enumerate}[leftmargin=1.5em]
  \item \textbf{长度偏置。} 原始 objective 里常见的 $1/|o_i|$ 归一化会影响长短回答的梯度权重，可能鼓励错误回答变长，或改变正负样本的更新强度。
  \item \textbf{无效 group。} 如果同一道题采样的一组回答全对或全错，组内 advantage 接近 $0$，这批样本几乎没有学习信号。
  \item \textbf{clip 太死。} 固定对称 clipping 可能过早压制有用探索，导致 entropy collapse。
  \item \textbf{token-level ratio 高方差。} reward 通常是整段回答级别，但 GRPO 常在 token 级别算 importance ratio，长序列里噪声会累积。
  \item \textbf{局部归一化偏置。} 只在同一个 prompt 的小 group 内归一化 reward，可能带来 question-level difficulty bias 或小样本估计不稳。
  \item \textbf{过度思考。} 只看最终答案 reward 时，模型可能学会越来越长的 reasoning，而不一定更高效。
\end{enumerate}

\begin{summarybox}
\textbf{一句话：} GRPO 的后续改进，大多不是推翻 group-relative RL，而是在修它的归一化、采样、clip、长度和粒度问题。
\end{summarybox}

\subsubsection{改进路线图}

\begin{center}
{\small
\begin{tabular}{p{2.6cm}p{4.4cm}p{5.8cm}}
\hline
\textbf{方法} & \textbf{主要修的问题} & \textbf{核心想法} \\
\hline
DAPO
&
无效 group、熵坍缩、长 CoT 训练不稳
&
动态采样、asymmetric clip、token-level loss、overlong reward shaping
\\
Dr.GRPO
&
GRPO 的长度偏置和 std 归一化偏置
&
去掉 $1/|o_i|$ 和 reward std scaling，用常数归一化恢复更无偏的目标
\\
GSPO
&
token-level importance ratio 高方差
&
把 importance ratio、clip、optimization 从 token level 移到 sequence level
\\
REINFORCE++ / RLOO
&
critic-free RL 的 advantage 估计
&
用全局 advantage normalization 或 leave-one-out baseline 替代 critic
\\
DCPO
&
固定 clipping 和相同 reward 组导致零梯度
&
动态 clipping、smooth advantage standardization，提高有效梯度利用率
\\
TreeRPO / WS-GRPO
&
只用最终答案 reward，过程信号弱或 rollout 过长
&
用树采样或弱监督信号给中间步骤 / prefix 更细的训练信号
\\
P-GRPO
&
不同用户 / 偏好分布混在一起
&
按 preference group 的历史 reward 归一化，避免主流偏好淹没少数偏好
\\
\hline
\end{tabular}
}
\end{center}

\subsubsection{第一类：长度偏置}

原始 GRPO 常见写法中，对每条回答有：

\[
\frac{1}{|o_i|}
\sum_{t=1}^{|o_i|}
l_{i,t}
\]

这意味着每条 sequence 的 token loss 会先被自己的长度平均。

问题是：当 advantage 为正或负时，长度平均会改变不同长度回答的更新强度。尤其在长 CoT reasoning 中，模型可能学到不理想的长度行为。

Dr.GRPO 的核心批评是：

\[
\frac{1}{|o_i|}
\text{ 会引入 response-level length bias}
\]

因此 Dr.GRPO 建议用全局常数归一化，例如最大生成长度 $L$：

\[
\mathcal{L}_\text{Dr.GRPO}
=
-
\frac{1}{LG}
\sum_{i=1}^{G}
\sum_{t=1}^{|o_i|}
l_{i,t}
\]

\begin{intubox}
\textbf{直觉：} 原始 GRPO 是“每条回答自己平均一下”；Dr.GRPO 改成“大家用同一个固定尺度归一化”，避免长短回答被不同规则加权。
\end{intubox}

\subsubsection{第二类：无效 Group}

GRPO 依赖同组比较。如果一个 group 的 reward 是：

\[
[0,0,0,0]
\]

或者：

\[
[1,1,1,1]
\]

那么组内相对 advantage 几乎没有区分度。

DAPO 的 dynamic sampling 就是为了解决这个问题：只保留有区分度的 group，例如：

\[
0
<
\left|
\{o_i \mid \operatorname{is\_correct}(o_i)\}
\right|
<
G
\]

也就是说，一个有效 group 至少要有一个正确回答，也至少要有一个错误回答。

\begin{summarybox}
\textbf{直觉：} GRPO 需要“同题内部对比”。如果一组答案全一样好或全一样差，就没什么可学；DAPO 让训练更集中在有好坏差异的样本上。
\end{summarybox}

\subsubsection{第三类：Clip 和熵坍缩}

标准 PPO / GRPO 通常使用对称 clip：

\[
\operatorname{clip}(\rho, 1-\epsilon, 1+\epsilon)
\]

DAPO 的 Clip-Higher 使用不对称边界：

\[
\operatorname{clip}
(\rho, 1-\epsilon_\text{low}, 1+\epsilon_\text{high})
\]

通常令：

\[
\epsilon_\text{high} > \epsilon_\text{low}
\]

这样可以对提高好样本概率给更大的上升空间，同时仍然控制下降或异常更新。

\begin{intubox}
\textbf{直觉：} 对称 clip 有时太保守，可能过早压掉探索；Clip-Higher 给正向学习更多空间，帮助维持 entropy。
\end{intubox}

DCPO 进一步把 clip 做成动态的：不同 token 根据先验概率或训练状态使用不同 clip 范围，目标是减少 token 被过早截断导致的零梯度。

\subsubsection{第四类：Token-Level vs Sequence-Level}

GRPO 常用 token-level ratio：

\[
\rho_{i,t}
=
\frac{
\pi_\theta(o_{i,t}\mid q,o_{i,<t})
}{
\pi_{\theta_\text{old}}(o_{i,t}\mid q,o_{i,<t})
}
\]

但 reasoning reward 通常是 sequence-level：

\[
r_i = R(q,o_i)
\]

这会产生一个粒度不匹配：

\[
\text{reward 是整段级别，importance ratio 却逐 token 算}
\]

GSPO 的核心想法是改用 sequence-level ratio：

\[
\rho_i^\text{seq}
=
\left[
\frac{
\pi_\theta(o_i \mid q)
}{
\pi_{\theta_\text{old}}(o_i \mid q)
}
\right]^{1/|o_i|}
\]

等价写成 log 形式：

\[
\rho_i^\text{seq}
=
\exp
\left(
\frac{1}{|o_i|}
\sum_{t=1}^{|o_i|}
\log
\frac{
\pi_\theta(o_{i,t}\mid q,o_{i,<t})
}{
\pi_{\theta_\text{old}}(o_{i,t}\mid q,o_{i,<t})
}
\right)
\]

\begin{summarybox}
\textbf{直觉：} GRPO 是每个 token 都单独算 ratio；GSPO 是整段回答算一个 sequence ratio。它试图让“优化单位”和“reward 单位”更一致。
\end{summarybox}

\subsubsection{第五类：局部归一化 vs 全局归一化}

标准 GRPO 对同一个 prompt 的 group 做归一化：

\[
\hat{A}_i
=
\frac{r_i-\mu_\text{group}}{\sigma_\text{group}+\epsilon}
\]

这叫 prompt-level / local normalization。

它的优点是每道题内部可比，缺点是 group 太小，估计可能不稳；而且不同题目的 reward 方差不同，可能引入 question-level difficulty bias。

REINFORCE++ 的一个方向是使用 global advantage normalization：

\[
\hat{A}_i
=
\frac{r_i-\mu_\text{batch}}{\sigma_\text{batch}+\epsilon}
\]

其中均值和方差来自更大的 global batch，而不是单个 prompt 的小 group。

\begin{intubox}
\textbf{直觉：} GRPO 是“每道题内部比”；REINFORCE++ 更像“在更大的 batch 里统一标尺”。这可能让 advantage 估计更稳，但也会弱化 prompt 内部对比的纯粹性。
\end{intubox}

\subsubsection{第六类：过度思考和过程信号}

RLVR 常用最终答案 reward：

\[
R(q,o)=
\mathbb{1}[\text{final answer correct}]
\]

这个信号简单、客观，但只告诉模型最后对不对，不告诉模型中间哪一步开始变好或变坏。

这可能带来：

\begin{itemize}[leftmargin=1.5em]
  \item 模型越写越长；
  \item 有些 token 只是冗余思考，但仍然被同一个正 reward 鼓励；
  \item 错误答案里的长推理也可能被复杂的归一化和长度项影响。
\end{itemize}

TreeRPO、WS-GRPO 这类方法试图给中间步骤或 partial trajectory 更细的信号：

\[
\text{最终 reward}
\to
\text{prefix / step-level guidance}
\]

\begin{summarybox}
\textbf{直觉：} 标准 GRPO 像“只看最终卷面分”；过程型改进希望知道“哪一步开始走对，哪一步继续写是浪费”。
\end{summarybox}

\subsubsection{接下来的学习顺序}

后面不用一下子全学。比较稳的顺序是：

\begin{enumerate}[leftmargin=1.5em]
  \item \textbf{DAPO：} 最实用，包含 dynamic sampling、Clip-Higher、token-level loss、overlong reward shaping；
  \item \textbf{Dr.GRPO：} 理解 GRPO / DAPO 的 length bias 和 std scaling bias；
  \item \textbf{GSPO：} 理解 token-level ratio 和 sequence-level ratio 的差异；
  \item \textbf{REINFORCE++ / RLOO：} 理解 critic-free RL 不一定要 group std normalization；
  \item \textbf{DCPO / TreeRPO / P-GRPO：} 作为扩展，分别看动态 clip、过程信号和个性化偏好。
\end{enumerate}

\begin{summarybox}
\textbf{本轮最小结论：} GRPO 后续改进可以按五个问题理解：
\[
\text{长度怎么处理？}
\]
\[
\text{group 有没有有效学习信号？}
\]
\[
\text{clip 会不会压死探索？}
\]
\[
\text{ratio 应该按 token 还是 sequence 算？}
\]
\[
\text{advantage 应该局部归一化还是全局归一化？}
\]
先学 DAPO，再学 Dr.GRPO 和 GSPO，是最自然的路线。
\end{summarybox}

\subsubsection{本轮检查题}

\begin{enumerate}[leftmargin=1.5em]
  \item 为什么全对或全错的 group 对 GRPO 学习信号很弱？
  \item DAPO 的 dynamic sampling 想解决什么问题？
  \item Dr.GRPO 为什么关注 $1/|o_i|$ 带来的长度偏置？
  \item GSPO 为什么要把 token-level ratio 改成 sequence-level ratio？
  \item local normalization 和 global normalization 的直觉差别是什么？
\end{enumerate}
